\chapter{基本方法与排版示例}
\label{chap:examples}
\section{公式与符号说明}
公式应作为论述的一部分出现，先介绍计算目的，再给出表达式。样本均值可写为
\begin{equation}
\bar{x}=\frac{1}{n}\sum_{i=1}^{n}x_i .
\label{eq:mean}
\end{equation}
\begin{formulaexplain}
\item[$x_i$] 第 $i$ 个观测值；
\item[$n$] 样本数量，且 $n>0$；
\item[$\bar{x}$] 样本均值。
\end{formulaexplain}
引用公式时使用标签，例如由\eqrefcn{eq:mean}可计算观测数据的平均水平。多行公式可使用对齐环境：
\begin{align}
 e_i &= x_i-\bar{x},\\
 s^2 &= \frac{1}{n-1}\sum_{i=1}^{n}e_i^2,\qquad n>1.
\label{eq:variance}
\end{align}
\begin{definition}[样本集合]
由 $n$ 个观测值组成的集合记为 $\mathcal{X}=\{x_1,\ldots,x_n\}$。
\end{definition}
\begin{theorem}[离均差性质]
对于按\eqrefcn{eq:mean}计算的均值，有 $\sum_{i=1}^{n}(x_i-\bar{x})=0$。
\end{theorem}
\begin{proof}
将均值定义代入求和式，得到 $\sum_{i=1}^{n}x_i-n\bar{x}=0$。
\end{proof}

\section{插图与双语图题}
图应尽量靠近正文首次引用的位置。\figrefcn{fig:workflow}是直接由 LaTeX 绘制的通用流程示意图，不依赖原论文图片。实际插入外部图片时可将文件放入 figures 文件夹，并使用插图命令。
\begin{figure}[htbp]
\centering
% 原创通用示意图；无需图像文件或额外绘图工具。
\begingroup
\setlength{\fboxsep}{8pt}
\fbox{\parbox[c][20pt][c]{0.22\linewidth}{\centering 问题定义}}
\hspace{0.01\linewidth}$\longrightarrow$\hspace{0.01\linewidth}
\fbox{\parbox[c][20pt][c]{0.22\linewidth}{\centering 方法设计}}
\hspace{0.01\linewidth}$\longrightarrow$\hspace{0.01\linewidth}
\fbox{\parbox[c][20pt][c]{0.22\linewidth}{\centering 验证分析}}
\endgroup

\sylufigcaption[fig:workflow]{研究流程示意图}{Illustration of a research workflow}
\end{figure}
正文应解释图中的要素及其联系，避免只写“结果如图所示”。图中的文字应清晰可辨，比例和标注方式应与正文保持一致。

\section{三线表与数据说明}
\tabrefcn{tab:demo}展示三线表和双语表题。表内数值仅用于演示排版，不对应任何真实方法的性能。正式论文应写明单位、统计方式和实验条件。
\begin{table}[htbp]
\sylutabcaption[tab:demo]{虚构数据排版示例}{Typesetting example with fictional data}
\begin{tabular}{lcc}
\toprule
方法 & 示例指标一 & 示例指标二（秒）\\
\midrule
方法甲 & 0.80 & 1.20\\
方法乙 & 0.85 & 1.35\\
方法丙 & 0.88 & 1.50\\
\bottomrule
\end{tabular}
\end{table}
\section{算法环境}
算法\ref{alg:mean}展示输入、输出和循环的排版。伪代码宜突出主要步骤，具体实现细节可以放在附录中。
\begin{algorithm}[htbp]
\caption{样本均值计算示例}\label{alg:mean}
\begin{algorithmic}[1]
\REQUIRE 观测序列 $x_1,\ldots,x_n$，其中 $n>0$
\ENSURE 样本均值 $\bar{x}$
\STATE $s\leftarrow 0$
\FOR{$i=1,\ldots,n$}
\STATE $s\leftarrow s+x_i$
\ENDFOR
\STATE $\bar{x}\leftarrow s/n$
\RETURN $\bar{x}$
\end{algorithmic}
\end{algorithm}
\section{本章小结}
本章展示常见学术内容的排版方式。复制示例时应保留标签与引用的对应关系，并将说明文字替换为实际研究内容。
